\structure{КОНСТРУКТОРСКАЯ~ЧАСТЬ}

В конструкторской части рассматриваются вопросы проектирования разрабатываемой
системы распределённых вычислений и методики верификации соответствия TLA+
спецификации и исходного кода. Сначала приводится обоснование выбора алгоритма
консенсуса, языка программирования и инструментов реализации, после чего
описывается архитектура системы распределённых вычислений, построенной с
использованием выбранного алгоритма.

Далее рассматриваются вопросы, относящиеся к системе верификации:
разрабатывается схема взаимодействия между её основными частями и
обосновывается решение ранее сформулированной математической постановки задачи.
Таким образом, в данной части определяется конструктивная основа как для
реализации распределённой системы, так и для построения механизма проверки её
соответствия формальной спецификации.
